% ===================================================================
%  TAULA N.º 9 — La Muntanya i el Balneari. — El Bosc i la Caça.
%  Traduction 2026 d'apres texto/io, controlee sur texto/fr et
%  texto/es. Consigne : voir texto/ca/10-taula-01.tex.
% ===================================================================

%%K t09-apar-1 apar
\VUsaut{4.0mm}
\VUtitre{15.0pt}{60}{TAULA N.º 9}
\VUsaut{3.0mm}

%%K t09-tit-1 sub
\VUcentre{12.0pt}{0}{\textit{Primera escena.}}
\VUsaut{3.0mm}

%%K t09-tit-2 sub
\VUpk{11.4pt}{40}{La Muntanya. --- La Ciutat balneària.}
\VUsaut{3.0mm}

%%K t09-c1-01-1 p
1. --- Fa vuit dies que sóc a Pratz. No puc expressar tota l'admiració que vaig
sentir quan em vaig trobar davant de la meravellosa \VUgras{serralada}
\textsuperscript{(1)} dels Pirineus, la \VUgras{cresta} \textsuperscript{(2)} de la
qual es retalla damunt el fons blau del cel com un festó gegantí.

%%K t09-c1-02-1 p
2. --- No m'imaginava que els \VUgras{cims} \textsuperscript{(3)} pirinencs
assolissin una \VUgras{altura} tan gran, i vaig quedar astorat davant els
magnífics \VUgras{glaceres} \textsuperscript{(5)} i els \VUgras{pics}
\textsuperscript{(4)} coberts de neu pels quals rodolen \VUgras{allaus} tan
terribles.

%%K t09-tit-3 sub
\VUsaut{3.0mm}
\VUpk{10.2pt}{40}{Paisatge de muntanya.}
\VUsaut{3.0mm}

%%K t09-c1-03-1 p
3. --- L'endemà mateix de la meva arribada vaig anar al vell \VUgras{volcà}
\textsuperscript{(6)} apagat que hi ha vora Pratz. Als vorals àrids i nus del
\VUgras{cràter} es veuen encara molt bé les petjades del pas de la \VUgras{lava}
que corria, fa uns quants segles, pel \VUgras{vessant de la muntanya}
\textsuperscript{(7)}. Vaig aconseguir trobar, en un dels \VUgras{pendents},
alguns fragments d'aquella lava.

%%K t09-c1-04-1 p
4. --- Pratz és a la frontera, i la \VUgras{fortalesa} \textsuperscript{(8)} que
defensa el \VUgras{coll} \textsuperscript{(27)} que separa França d'Espanya
recorda del tot aquells vells castells de les vores del Rin que semblen aguaitar
una \VUgras{presa} des de dalt de la seva roca.

%%K t09-c1-05-1 p
5. --- Una \VUgras{àguila} \textsuperscript{(9)} enorme, que té el seu niu no lluny
del \VUgras{fort} \textsuperscript{(8)}, atrau sovint la meva atenció. Aquest
\VUgras{ocell de rapinya}, de \VUgras{bec} poderós, porta sovint a les seves
cries un anyell o un cabrit subjectat a les seves \VUgras{urpes}. Les seves
\VUgras{ales} són tan grans que pot planar molta estona amb la seva pesada
càrrega.

%%K t09-tit-4 sub
\VUsaut{3.0mm}
\VUpk{10.2pt}{40}{L'Excursionista.}
\VUsaut{3.0mm}

%%K t09-c1-06-1 p
6. --- El passeig més agradable d'aquí és l'excursió a la \VUgras{cascada}
\textsuperscript{(10)} de Cau. El \VUgras{salt d'aigua} cau en remolins i després
es perd en un bonic \VUgras{rierol} a la \VUgras{pineda d'avets}
\textsuperscript{(11)} situada a l'oest de la ciutat. Vam pujar per aquell bosc
fins a l'\VUgras{observatori meteorològic} \textsuperscript{(12)} i, des de dalt
del \VUgras{mirador}, vam abastar el \VUgras{panorama} sencer de la comarca. El
\VUgras{paisatge} és superb i la \VUgras{natura} sembla haver prodigat a aquest
país tots els tresors de què disposa.

%%K t09-c1-07-1 p
7. --- Per baixar vam agafar el \VUgras{funicular} \textsuperscript{(13)} i ens vam
aturar en una petita \VUgras{estació} vora les \VUgras{ruïnes}
\textsuperscript{(14)} d'un vell \VUgras{castell feudal} habitat antany pels
senyors de Pratz.

%%K t09-c1-08-1 p
8. --- Pobre castell! Què deu pensar en veure als seus peus la coqueta
\VUgras{ciutat balneària} \textsuperscript{(15)}, que avui és una \VUgras{estació
termal} de moda?

%%K t09-c1-08-2 p
El \VUgras{clima} és tan agradable i l'\VUgras{aigua mineral} tan benèfica que la
\VUgras{cura} que se segueix al \VUgras{sanatori} \textsuperscript{(17)} és un
veritable plaer.

%%K t09-tit-5 sub
\VUsaut{3.0mm}
\VUpk{10.2pt}{40}{Una agradable ciutat termal.}
\VUsaut{3.0mm}

%%K t09-c1-09-1 p
9. --- Altrament, s'ha fet tot perquè l'estada sigui agradable. Es va construir un
gran \VUgras{viaducte} \textsuperscript{(16)} per poder estendre-hi una via fèrria;
el \VUgras{parc} \textsuperscript{(18)} del \VUgras{balneari}, on es donen
\VUgras{concerts}, està arranjat d'una manera encantadora. S'han construït
diversos graciosos \VUgras{xalets} \textsuperscript{(19)} al voltant del
\VUgras{casino} \textsuperscript{(21)}, i una \VUgras{calçada}
\textsuperscript{(22)} molt ben cuidada fa facilíssimes les comunicacions.

%%K t09-c1-10-1 p
10. --- Un simple \VUgras{talús} \textsuperscript{(23)} separa la \VUgras{carretera}
\textsuperscript{(20)} de la \VUgras{vall} \textsuperscript{(25)} pròpiament dita,
al fons de la qual s'estén un bonic \VUgras{llac} \textsuperscript{(26)} que
alimenta un \VUgras{rierol} \textsuperscript{(24)} molt clar.

%%K t09-c1-11-1 p
11. --- A l'altre costat de la vall s'explota una \VUgras{mina}
\textsuperscript{(28)} de ferro. Vaig anar diverses vegades a veure baixar els
\VUgras{miners} pel \VUgras{pou}, i fins i tot vaig entrar a la \VUgras{galeria}
on passen el dia extraient el \VUgras{mineral} que conté el \VUgras{metall}.

%%K t09-c1-12-1 p
12. --- Es va construir una important \VUgras{foneria} vora la mina per aprofitar
de seguida les riqueses que se'n treuen. L'\VUgras{alt forn}
\textsuperscript{(29)}, escalfat amb el \VUgras{carbó de pedra} que abunda aquí,
permet d'obtenir \VUgras{ferro colat} d'una manera ràpida i barata.

%%K t09-c1-13-1 p
13. --- No lluny de la mina s'obre una \VUgras{pedrera} \textsuperscript{(30)}. No
és ni \VUgras{granit}, ni \VUgras{pòrfir}, ni tan sols \VUgras{gres} el que el
vell \VUgras{picapedrer} arrenca amb el seu \VUgras{pic}, sinó simple
\VUgras{pedra de carreu} per construir cases.

%%K t09-c1-14-1 p
14. --- Un dels més bells ornaments d'aquest país és l'\VUgras{avet}
\textsuperscript{(31)}. Per tot arreu --- al fons d'un \VUgras{barranc}
\textsuperscript{(32)}, a la vora d'un \VUgras{torrent} \textsuperscript{(33)},
vora un \VUgras{avenc} \textsuperscript{(34)} o un precipici ---, les seves fines
agulles hi posen la seva nota verda.

%%K t09-tit-6 sub
\VUsaut{3.0mm}
\VUpk{10.2pt}{40}{Caminar a peu.}
\VUsaut{3.0mm}

%%K t09-c1-15-1 p
15. --- Aprofito les meves vacances per fer llargs passeigs. Els \VUgras{camins}
\textsuperscript{(35)} que serpentegen per la muntanya permeten de recórrer, sense
adonar-se de la distància, diversos \VUgras{quilòmetres} i sovint diverses
\VUgras{llegües}.

%%K t09-c1-15-2 p
\VUgras{Fites} \textsuperscript{(36)} i \VUgras{indicadors}
\textsuperscript{(37)} jalonen els camins, on hom troba sovint un jove
\VUgras{muntanyenc} \textsuperscript{(38)} amb el \VUgras{sarró}
\textsuperscript{(40)} al costat, que porta una \VUgras{marmota}
\textsuperscript{(39)}, o algun \VUgras{gitano} \textsuperscript{(41)} la
\VUgras{gorra de pell} \textsuperscript{(42)} del qual contrasta d'una manera
singular amb el seu vell \VUgras{caputxó} \textsuperscript{(43)} trencat. Porta de
la \VUgras{cadena} un \VUgras{ós} \textsuperscript{(44)} ensinistrat.

%%K t09-tit-7 sub
\VUpk{10.2pt}{40}{Escalar muntanyes.}
\VUsaut{3.0mm}

%%K t09-c1-16-1 p
16. --- Gairebé cada dia vaig amb un \VUgras{guia} \textsuperscript{(48)} a fer una
\VUgras{ascensió}, i estic molt orgullós quan, amb la \VUgras{motxilla}
\textsuperscript{(47)} a l'esquena i el \VUgras{bastó alpí} a la mà, com un
veritable \VUgras{turista} \textsuperscript{(46)}, escalo les \VUgras{roques}
\textsuperscript{(45)}.

%%K t09-c1-17-1 p
17. --- Des de dalt d'una muntanya, els habitants de la plana no semblen més
grans que formigues; un \VUgras{ramat d'ovelles} \textsuperscript{(50)} que
pastura en una \VUgras{pastura} \textsuperscript{(49)} sembla una joguina de nen.
Un \VUgras{pi} \textsuperscript{(51)}, o fins i tot una \VUgras{cabana de fusta}
\textsuperscript{(52)}, amb prou feines són una taca en la verdor.

%%K t09-c1-18-1 p
18. --- Durant aquestes excursions ens creuem sovint al \VUgras{sender}
\textsuperscript{(53)} amb un \VUgras{caçador d'isards} \textsuperscript{(54)} que
porta a l'espatlla l'animal que acaba de matar.

%%K t09-c1-19-1 p
19. --- He posat al meu herbari una branca de \VUgras{falguera}
\textsuperscript{(55)} molt curiosa i un bonic ramet rosat de \VUgras{bruc}
\textsuperscript{(57)}, que vaig collir a l'entrada d'una petita \VUgras{cova}
\textsuperscript{(56)} en el meu darrer passeig.

%%K t09-tit-8 sub
\VUsaut{3.0mm}
\VUcentre{12.0pt}{0}{\textit{Segona escena.}}
\VUsaut{3.0mm}

%%K t09-tit-9 sub
\VUpk{11.4pt}{40}{El Bosc. --- La Caça.}
\VUsaut{3.0mm}

%%K t09-c2-01-1 p
1. --- Ahir vaig passar un dia deliciós al bosc. L'activitat que regna entre els
animals i els \VUgras{insectes} \textsuperscript{(5)} que el poblen em va
interessar vivament.

%%K t09-c2-01-2 p
L'\VUgras{aranya} \textsuperscript{(1)} que teixeix la seva \VUgras{teranyina}
\textsuperscript{(2)} entre dues branques per atrapar la \VUgras{mosca}
\textsuperscript{(3)} que passa, el \VUgras{cargol} \textsuperscript{(4)} que
porta penosament la seva closca, el cérvol volant que busca la seva presa, la
formiga diligent que s'afanya vora el \VUgras{formiguer} \textsuperscript{(6)}:
tot aquell petit món anava, venia i treballava amb una animació extraordinària.

%%K t09-c2-02-1 p
2. --- Una \VUgras{serp} \textsuperscript{(7)}, que al principi vaig prendre per un
\VUgras{escurçó} però que no era més que una simple \VUgras{colobra}, estava
arraulida sota un tronc d'arbre i de tant en tant treia el seu caparró fi, que
sempre semblava disposat a mossegar. Davant meu, un gros \VUgras{llimac}
\textsuperscript{(8)} reptava pesadament després de devorar una de les saboroses
\VUgras{maduixes de bosc} \textsuperscript{(11)} que hi creixen en abundància.

%%K t09-c2-03-1 p
3. --- El sòl era encatifat de \VUgras{flors del bosc}: \VUgras{prímules}
\textsuperscript{(9)}, \VUgras{vincapervinques} \textsuperscript{(10)} i moltes
altres plantes que jo no coneixia florien entre les \VUgras{mates d'herba}
\textsuperscript{(12)}.

%%K t09-c2-04-1 p
4. --- En aquesta comarca, la \VUgras{tòfona} és gairebé tan comuna com el
\VUgras{bolet} \textsuperscript{(14)}, i un \VUgras{porcell}
\textsuperscript{(13)} d'un guarda forestal buscava amb aire golut a veure si en
trobava alguna.

%%K t09-tit-10 sub
\VUsaut{3.0mm}
\VUpk{10.2pt}{40}{La caça furtiva.}
\VUsaut{3.0mm}

%%K t09-c2-05-1 p
5. --- Vaig assistir al \VUgras{final} d'una escena que em va divertir molt.
Gràcies a l'olfacte del seu \VUgras{gos bàsset} \textsuperscript{(15)}, el
\VUgras{guarda forestal} \textsuperscript{(16)} acabava de sorprendre un
\VUgras{caçador furtiu} \textsuperscript{(22)} en el moment en què aquest es
disposava a endur-se la \VUgras{caça} \textsuperscript{(17)} que havia matat.
Preferint abandonar la seva presa a deixar-se detenir, el furtiu havia fugit, i
\VUgras{llebre} \textsuperscript{(18)}, \VUgras{cérvola}
\textsuperscript{(19)}, \VUgras{cabirol} \textsuperscript{(20)} i
\VUgras{senglar} \textsuperscript{(21)} quedaven damunt l'herba, per alegria del
guarda.

%%K t09-c2-06-1 p
6. --- La varietat d'arbres que tanca el bosc és sorprenent. Els \VUgras{faigs}
\textsuperscript{(23)} i els \VUgras{bedolls} \textsuperscript{(26)}, els
\VUgras{pins}, que donen la \VUgras{resina} \textsuperscript{(27)}, els
\VUgras{roures} \textsuperscript{(29)} i molts més entremesclen les seves
branques.

%%K t09-c2-06-2 p
L'\VUgras{heura} \textsuperscript{(24)}, que s'agafa al tronc dels arbres alts, dóna
a la \VUgras{devesa} \textsuperscript{(25)} un verd brillant que s'uneix
agradablement al to més clar i més suau de la \VUgras{molsa}
\textsuperscript{(31)} en què el \VUgras{picot verd} \textsuperscript{(30)} busca
insectes.

%%K t09-tit-11 sub
\VUsaut{3.0mm}
\VUpk{10.2pt}{40}{Paisatge de bosc.}
\VUsaut{3.0mm}

%%K t09-c2-07-1 p
7. --- Els vells roures em van encantar. Les seves gruixudes \VUgras{arrels}
\textsuperscript{(32)} recargolades, mig sortides de la terra, els donen una
esplèndida aparença de força i de vigor, i em pregunto com el
\VUgras{propietari} \textsuperscript{(35)} pot resignar-se a fer-los talar i a
lliurar-los a la \VUgras{serra} \textsuperscript{(34)} del \VUgras{llenyataire}
\textsuperscript{(33)}. Els pobres \VUgras{troncs} \textsuperscript{(36)},
despullats de la seva \VUgras{escorça} \textsuperscript{(37)} o esberlats per la
\VUgras{destral} \textsuperscript{(38)} del \VUgras{jornaler}
\textsuperscript{(39)}, em fan pena.

%%K t09-c2-08-1 p
8. --- A prop d'allà, un vell \VUgras{carboner} \textsuperscript{(41)} havia
preparat amb la seva \VUgras{pala} una \VUgras{pila} \textsuperscript{(42)} de
carbó, i una vella s'enduia un \VUgras{feix} \textsuperscript{(40)} de
\VUgras{llenya morta} fet amb els branquillons dels arbres talats.

%%K t09-tit-12 sub
\VUsaut{3.0mm}
\VUpk{10.2pt}{40}{La Cacera.}
\VUsaut{3.0mm}

%%K t09-c2-09-1 p
9. --- Volia anar a descansar una estona a la \VUgras{casa del guarda forestal}
\textsuperscript{(46)}, però quan vaig arribar vora l'\VUgras{estany}
\textsuperscript{(43)}, en què neda un bell \VUgras{cigne} \textsuperscript{(45)},
em vaig adonar que una \VUgras{inundació} \textsuperscript{(44)} recent havia
convertit el camí en un veritable \VUgras{aiguamoll}. Vaig haver de fer una volta
i, en passar vora el \VUgras{cedre} \textsuperscript{(49)}, els \VUgras{pollancres}
\textsuperscript{(48)} i l'\VUgras{om} \textsuperscript{(47)} que són a la
llinda del bosc, vaig veure sortir d'un \VUgras{matoll} \textsuperscript{(54)} un
magnífic \VUgras{cérvol} \textsuperscript{(50)} perseguit per una \VUgras{gossada}
\textsuperscript{(51)} enfurismada, a la qual excitava a més el \VUgras{corn}
\textsuperscript{(53)} dels \VUgras{caçadors a cavall} \textsuperscript{(52)}. El
pobre animal estava exhaust. Va venir a caure moribund a uns passos de mi.

%%K t09-c2-10-1 p
10. --- El fill del guarda forestal, atret pel soroll de la \VUgras{cacera},
va sortir de la casa i vam tornar tots dos al bosc. Havia vist una
\VUgras{garsa} \textsuperscript{(56)} a la branca d'un vell roure. Pensava que el
seu niu devia estar amagat al \VUgras{fullatge} \textsuperscript{(58)} i el volia
agafar.

%%K t09-c2-10-2 p
Àgil com un \VUgras{esquirol} \textsuperscript{(61)}, va enfilar-se de
\VUgras{branca} \textsuperscript{(55)} en branca, però no va trobar res i només va
aconseguir fer volar una vella \VUgras{cornella} \textsuperscript{(60)} posada en
un \VUgras{branquilló gruixut} \textsuperscript{(57)}.
\VUsaut{4.0mm}
\VUfilet{22mm}
